\documentclass[10pt,a4paper]{amsart}
\usepackage[margin=19mm]{geometry}
\usepackage{amsmath,amssymb,mathtools}
\usepackage[scr=rsfs,cal=euler]{mathalfa}
\usepackage{xfrac}
\usepackage[unicode,hidelinks,bookmarks=false]{hyperref}
\newcommand{\ie}{i.e.\ }
\newcommand{\cf}{cf.\ }
\newcommand{\resp}{respectively\ }
\newcommand{\Subsection}{\subsection}
\providecommand{\nobreakdash}{\nobreak\hbox{-}}
\setlength{\emergencystretch}{2em}
\multlinegap0pt
\usepackage{bm}
\usepackage{bbm}
\usepackage{dsfont}
\usepackage{xspace}
\usepackage{cancel}
\usepackage{mathtools}
\usepackage{amsthm}
\makeatletter
\@ifundefined{theorem}{\newtheorem{theorem}{Theorem}}{}
\@ifundefined{lemma}{\newtheorem{lemma}[theorem]{Lemma}}{}
\makeatother
\title[On Weighted and Bounded Multidimensional Catalan Numbers]{On Weighted and Bounded Multidimensional Catalan Numbers}
\author{Ryota Inagaki, Dimana Pramatarova}
\date{Source: arXiv:2510.14956v1 (CC BY 4.0)}
\begin{document}
\maketitle

\noindent\emph{Source and licence.} On Weighted and Bounded Multidimensional Catalan Numbers. Version arXiv:2510.14956v1 (CC BY 4.0), distributed under the Creative Commons Attribution 4.0 International licence, as recorded in the arXiv licence metadata for this exact version and shown on the abstract page https://arxiv.org/abs/2510.14956v1. Full rights notice: licenses/arxiv-2510.14956-CC-BY-4.0.txt.\par\smallskip

\noindent\textit{Editorial scope.} Counted unit: the Theorem whose source label is \texttt{theorem:general-periodicity} in the article (source0.tex lines 539--542, source label \texttt{theorem:general-periodicity}), delivered together with the article's own complete original proof of it (source0.tex lines 544--570, one \texttt{proof} environment of 27 lines). No mathematics is written, repaired or re-derived: statement and proof are the article's own text line for line. Required context retained uncounted: lemma:transmatrixcat (lines 507--529). Every definition, display and label of those lines is retained unchanged. Omitted from this package and not redistributed: 1--506, 530--538, 543--543, 571--1072. Nothing inside a retained excerpt is omitted or altered: every retained body is delivered exactly as the article's lines are written, so no omission receipt of any kind accompanies this package. Citations: the retained excerpts carry no citation command at all, so no bibliography entry of the article is transcribed and no reference is redistributed. Reference adaptation: none was needed and none was made; every reference inside the delivered unit resolves inside it, so no unresolved reference or citation is produced. Printed slips kept literal: no printed word, symbol, hypothesis or number of the counted statement or of its proof was altered. Layout and class: the article's own class is \texttt{amsart} with packages including graphicx, amsmath, amsfonts, amssymb, amsthm, fullpage, inputenc, geometry, enumitem, xcolor, caption, tabularx, fancyhdr, comment; the wrapper is the lane's pinned \texttt{amsart} editorial wrapper at 10pt on A4, which restates the theorem-like environments the retained text needs and adds the article's own macro definitions that the retained text needs (none), with extra wrapper packages (bm, bbm, dsfont, xspace, cancel, mathtools). The delivered numbering is the unit's own. Assets: no retained span contains a figure environment, a graphics-inclusion command, a PDF-page inclusion, a table or a TikZ picture; no image file is copied into this package, the package declares no asset dependency and no image file is redistributed with it. Licence: version arXiv:2510.14956v1 of this article is distributed under the Creative Commons Attribution 4.0 International licence, as recorded in the arXiv licence metadata for this exact version and shown on the abstract page https://arxiv.org/abs/2510.14956v1. A full rights notice is delivered at licenses/arxiv-2510.14956-CC-BY-4.0.txt. Characters: every retained excerpt is pure ASCII, so the delivered engine, XeLaTeX with its default Latin Modern text font, reports no missing character.
\begin{lemma} The $2$-dimensional weighted Catalan numbers satisfy the following recurrence:
\label{lemma:transmatrixcat}
 %Denote by $C^{\vec{b}}_{n,i}$ the total weight of the Dyck paths from $(0,i)$ to $(2n,0)$. $($In particular, $C^{\vec{b}}_{n,0} = C^{\vec{b}}_{n})$.  Then
 \[\begin{bmatrix}
C^{\vec{b}}_{n,0} \\
C^{\vec{b}}_{n,2} \\
C^{\vec{b}}_{n,4} \\
\vdots
\end{bmatrix} = 
\begin{bmatrix}
b_0 & b_0 b_1 & 0 & \dots  &\dots &\dots \\
1 & b_1 + b_2 & b_2 b_3 & 0 & \dots &\dots \\
0 & 1 & b_3 + b_4 & b_4 b_5 & 0 & \dots \\
\vdots & \vdots & \vdots & \vdots & \ddots &\dots
\end{bmatrix}
\begin{bmatrix}
C^{\vec{b}}_{n-1,0} \\
C^{\vec{b}}_{n-1,2} \\
C^{\vec{b}}_{n-1,4} \\
\vdots
\end{bmatrix}.
\] 
\end{lemma}

\begin{theorem}
    \label{theorem:general-periodicity}
    For any positive integer $m$, the sequence $C^{\vec{b}}_{1}, C^{\vec{b}}_2, \dots$ is periodic modulo $m$ if $m$ divides $b_0b_1\ldots b_k$ for some non-negative integer $k$. 
\end{theorem}

\begin{proof}
Recall that $C^{\vec{b}}_{n,i}$ is the number of weighted Dyck paths from $(0,i)$ to $(2n,0)$. Observe that the weight of each Dyck path for $i > \left\lfloor \frac{k}{2} \right\rfloor$ is divisible by $b_0b_1\ldots b_k$. Together with Lemma \ref{lemma:transmatrixcat}, this implies that the transition matrix modulo $m$ is of finite size $(\ell + 1) \times (\ell + 1)$, which depends on the parity of $k$. Specifically, $\ell = \lfloor \frac{k}{2} \rfloor.$ If $k$ is even, then the last element of the matrix is $a_{2\ell} = b_{k-1}$, and if $k$ is odd, then it is ${a_{2\ell}} = b_{k-2}+b_{k-1}$. 

\[\begin{bmatrix}
C^{\vec{b}}_{n,0} \\
C^{\vec{b}}_{n,2} \\
C^{\vec{b}}_{n,4} \\
\vdots \\
C^{\vec{b}}_{n,2 \ell}
\end{bmatrix} = 
\begin{bmatrix}
b_0 & b_0 b_1 & 0 & \dots  &\dots &\dots \\
1 & b_1 + b_2 & b_2 b_3 & 0 & \dots &\dots \\
0 & 1 & b_3 + b_4 & b_4 b_5 & 0 & \dots \\
\vdots & \vdots & \vdots & \vdots & \ddots &\dots \\
0 & 0 & \dots & 0 & 1 & a_{2 \ell}
\end{bmatrix}
\begin{bmatrix}
C^{\vec{b}}_{n-1,0} \\
C^{\vec{b}}_{n-1,2} \\
C^{\vec{b}}_{n-1,4} \\
\vdots \\
C^{\vec{b}}_{n-1,2 \ell}
\end{bmatrix}.
\]

There exist positive integers $s$ and $t$ such that $(C^{\vec{b}}_{s,0}, \ldots, C^{\vec{b}}_{s,2\ell}) \equiv  (C^{\vec{b}}_{s+t,0}, \ldots, C^{\vec{b}}_{s+t,2\ell})$, because by the Pigeonhole principle, there are at most $m^{\ell+1}$ possible combinations of $(C^{\vec{b}}_{n,0}, \ldots, C^{\vec{b}}_{n,2\ell}) \pmod m$. The matrix is of finite size, and thus the sequence will eventually be periodic, with $C^{\vec{b}}_{n+jt,0} \equiv C^{\vec{b}}_{n,0} \pmod m$ for any positive integer $j$. \end{proof}

\end{document}
